\section{Results}
\label{sec:results}

% Main results table (tab:transfer) is declared in signal.tex so it
% floats to an early page.

\subsection{Live escalation}
\label{sec:live}
\label{sec:liveloop}
\label{sec:livecurve}

We ran the 240 frozen test queries through four escalation tiers,
each a complete turn-based loop with audio input
(\S\ref{sec:liveproto}). The primary sweep retains the model's output
as text; a separate run of the balanced tier with speech synthesis
enabled leaves all 84/240 gate decisions unchanged and measures
speech onset. We score deployed-channel answer accuracy: the
correctness of the relayed content after the full system path.
Headline numbers use a pre-registered input filter that removes 21
formula-LaTeX queries and one broken recording, because the speech
rendering destroys the symbolic input before the model receives it.
Accuracy rises monotonically,
$\mathbf{42.2\% \to 52.8\% \to 59.2\% \to 63.3\%}$ at 0/14/31/57\%
escalation, and every tier's gain is significant (paired bootstrap;
Table~\ref{tab:transfer}, Figure~\ref{fig:dualview}). The gate also
exceeds matched-rate random
escalation at every tier: by 7.2--9.4 points on deployed-channel
accuracy, and by 5.3--8.3 points when the relay channel is factored
out (Figure~\ref{fig:dualview}).

% fig:dualview is declared at the end of Section 4 so it floats
% onto the page where these results are discussed.

A measured always-escalate arm reaches 66.5\%. The gain from 57\% to
100\% escalation is small because the expert sees the model's own
transcription: the same sessions reach 92.2\% when the expert
receives the reference transcript, so the transcription channel, not
selection, is the limit at high escalation rates. Live accuracies
have a replication floor of $\pm$2--3 points; no claim rests on a
smaller difference. Full tables, latency profiles, and boundary cases
are in Appendix~\ref{app:livedetail}.

\paragraph{Serving-regime checks.}
Replaying the loop with the talker enabled leaves the frozen ranking
nearly unchanged (AUC 0.871/0.856 versus 0.866 with the talker
disabled). A teacher-forced concurrent-prefill variant is less stable
and requires an in-regime refit; its external results remain
exploratory (Appendix~\ref{app:duplexval}). Under the model's native
full-duplex head, the read moves to the head's own listen$\to$speak
commit point, and the probe is refit in regime on a widened
5{,}228-row calibration merge under the official serving
configuration, with labels judged on the native answers. The refit
reads at AUC 0.876 internally and 0.771 on the four English external
QA pools. The Mandarin reasoning pool is lower, at 0.621; a
1{,}500-question held-out Mandarin set of facts, knowledge, and
commonsense reads at 0.811, so the gap is multi-step reasoning
rather than language (Appendix~\ref{app:native}). Re-mixing the
native gate's decisions with cached expert outcomes beats
matched-rate random on all five English pools at both deployed
tiers; with per-language thresholds the Mandarin pool escalates at
the nominal rates and beats matched-rate random at the two lower
tiers. A full native live run on the internal 240 queries, under the
earlier inherited serving configuration, improves accuracy from
0.371 to 0.537 at 45\% escalation, below the prediction of 0.596
from the same re-mix arithmetic, because relaying the returned
answer is lossy. These checks support the ranking beyond the primary
turn-based loop, but they also show that read points, probe weights,
and thresholds must be calibrated for the serving regime
(Appendix~\ref{app:native}).
